\section{Full proofs}
\label{app:proofs}

\subsection{Proof of Theorem~\ref{thm:forward-forgetting}}
\label{app:proof-thm-ff}

Let $D$ and $D' = D \setminus \{D_{i^\star}\}$ be neighbouring datasets differing in
client $i^\star$'s data, and let
$\Theta^{(r_{\mathrm{req}})}_\mathrm{sh}(D)$ denote the feature-block output of
Algorithm~\ref{alg:decayed-aggregation} after $r_{\mathrm{req}}$ rounds on $D$.

\paragraph{Linear form of the aggregation state.}
From~\eqref{eq:feature-update},
$\theta^{(r_{\mathrm{req}})}_\mathrm{sh}
 = \theta^{(0)}_\mathrm{sh}
 + \eta_g \sum_{t=0}^{r_{\mathrm{req}}-1} \tilde\Delta^{(t+1)}_\mathrm{sh}
 + \eta_g \sum_{t=0}^{r_{\mathrm{req}}-1} \xi^{(t)}$,
where the second term is deterministic given the client updates and the third is a
sum of independent Gaussian variates with total covariance
$r_{\mathrm{req}} \eta_g^2 \sigma^2 I$.

\paragraph{Sensitivity to removing $i^\star$.}
Unrolling the EMA~\eqref{eq:decay-state} gives
$\tilde\Delta^{(t+1)}_\mathrm{sh} = (1-\alpha) \sum_{s=0}^{t} \alpha^{t-s}\,
\bar\Delta^{(s)}_\mathrm{sh}$, where each $\bar\Delta^{(s)}_\mathrm{sh}$ changes
in L2 by at most $C/|S^{(s)}|$ when $i^\star \in S^{(s)}$. Summing and bounding the
geometric series by $1/(1-\alpha)$, with $S \subseteq [0, r_\forget]$ and
$r_\mathrm{req} - s \ge T$ for all $s \in S$, yields
$\|\Delta_{i^\star}\|_2 \le \eta_g C \alpha^T$ under the participation bound
$|S| \le \rho r_\forget$.

\paragraph{Gaussian mechanism on the final release.}
The final state is a deterministic function of the data plus independent Gaussian
noise with per-coordinate standard deviation $\tau = \sqrt{r_{\mathrm{req}}} \eta_g
\sigma$. The analytic Gaussian mechanism (Theorem~A.1 of \cite{dwork2014book})
gives, for sensitivity $\|\Delta_{i^\star}\|_2$,
$\varepsilon_1(T) \le \|\Delta_{i^\star}\|_2 \sqrt{2\ln(1.25/\delta)} / \tau
= C \alpha^T \sqrt{2\ln(1.25/\delta)} / (\sigma \sqrt{r_\mathrm{req}})$, matching
Eq.~\eqref{eq:eps-rate}. \qed

\begin{remark}[Trajectory release]
The bound above applies to the released post-erasure state. For an adversary that
observes every intermediate round broadcast, Appendix~\ref{app:rdp} gives a
R\'enyi-DP accounting that retains the $\alpha^T$ decay and adds a $\sqrt{r_\mathrm{req}}$
trajectory factor.
\end{remark}

\subsection{Proof of Lemma~\ref{lem:subtraction}}
\label{app:proof-lem-sub}

Fix the training trajectory and $\theta_\mathrm{cl}$ at erasure time. Under
Algorithm~\ref{alg:decayed-aggregation}, $\theta_\mathrm{cl}$ is deterministic
given the sequence of client gradients and does not depend on any client summary.
The release $\sigma_{i^\star} = \mathrm{SS}(D_{i^\star}) + \eta_\summary$ is
$(\varepsilon_\summary, \delta_\summary)$-DP against $D_{i^\star}$ by construction
of the Gaussian mechanism. The erasure map $\mathcal{E}(\theta, \cdot)$ is a
deterministic post-processing of $\sigma_{i^\star}$, so by DP post-processing
(Proposition~2.1 of \cite{dwork2014book}),
$\theta'_\mathrm{cl} \indist[(\varepsilon_\summary, \delta_\summary)]
\theta'_\mathrm{cl, -i^\star}$. \qed

\subsection{Non-convex trajectory argument}
\label{app:nonconvex}

The sensitivity argument of Theorem~\ref{thm:forward-forgetting} uses only the
linear EMA structure. The non-linear dependence of local-training outputs
$\Delta_i^{(t)}$ on $\theta^{(t)}$ can be handled with a trajectory-Lipschitz
contraction following \cite{gupta2021adaptive}: under bounded local updates
(Assumption~\ref{ass:clipping}) and Lipschitz continuity of the local-training map
in its input state, the state deviation $\|\theta^{(t+1)}(D) - \theta^{(t+1)}(D')\|_2$
contracts geometrically whenever $\eta_g (1-\alpha) \eta_l L' < 1$, where $L'$ is
the per-round Lipschitz constant of the local-training map. This replaces $C$ by a
trajectory-adjusted constant $C'$ depending on $\eta_l, L', \alpha$; the $\alpha^T$
dependence is preserved.

\paragraph{Empirical $L'$ at ResNet-18 scale.}
We measured $L'$ directly on the paper's CIFAR-10 operating point. Taking the final
global state $\theta^{(200)}$ as the probe origin, we drew $4$ random unit
directions $v$ (restricted to trainable parameters, i.e., excluding BatchNorm
running statistics) and, for each of $8$ randomly sampled clients, ran local
training (3 epochs, batch 32, lr 0.02) from $\theta^{(200)}$ and from
$\theta^{(200)} + \varepsilon v$ ($\varepsilon = 10^{-3}$) with identical data
shuffling. The estimate $\hat L' = \| \mathrm{LocalTrain}(\theta^{(200)} + \varepsilon v) - \mathrm{LocalTrain}(\theta^{(200)}) \|_2 / \varepsilon$
across $32$ (client, direction) pairs gave
median $8.0$, mean $22.9$, p95 $87.0$, max $88.3$. All $32$ measurements fall
below the contraction threshold $L' < 100$ implied by the paper's hyperparameters
($\eta_g = 1$, $\alpha = 0.5$, $\eta_l = 0.02$), so the one-round contraction
factor $\eta_g (1-\alpha) \eta_l L'$ is at most $0.883$ at the operating point.
This gives an explicit numeric for the trajectory-adjusted constant: $C' \le C /
(1 - 0.883) \approx 8.5\,C$ in the worst case over our probe. The $\alpha^T$
decay dominates this correction factor for any $T \gtrsim 20$, so the
operating-point numerics in Theorem~\ref{thm:composition} are unaffected to within
a $9\times$ constant. The probe script is released with the code artefact.

\subsection{R\'enyi-DP tightening for full-trajectory release}
\label{app:rdp}

Under the full-trajectory release model, the per-round Gaussian mechanism on
$\tilde\Delta^{(t+1)}$ has sensitivity $\mathrm{sens}_t = \eta_g (1-\alpha) C
\alpha^{r_\mathrm{req}-t-1}$ and induces a per-round R\'enyi-DP guarantee of order
$\lambda$ equal to $\lambda\, \mathrm{sens}_t^2 / (2\sigma^2)$
\cite{mironov2017renyi}. Summing over $t \in S$ and converting to $(\varepsilon,
\delta)$-DP via the standard RDP conversion gives, for $|S| \le \rho r_\mathrm{req}$,
$\varepsilon \le \eta_g (1-\alpha) C \alpha^T \sqrt{2 \rho r_\mathrm{req} \ln(1/\delta)}
/ (\sigma \sqrt{1 - \alpha^2})$. The $\alpha^T$ decay is preserved; the trajectory
length contributes a $\sqrt{r_\mathrm{req}}$ factor.
